\documentclass[a4paper,11pt]{article}

\usepackage{fullpage}
\usepackage[all]{xy}
\usepackage[latin1]{inputenc}
\usepackage{amsmath,amsthm}
\usepackage{amsfonts}
%\usepackage{algorithm}
\usepackage{graphicx}
%\usepackage[compatible,noend]{algpseudocode}
\usepackage{latexsym}
\usepackage{float}
\usepackage{color}

%%%%%%%%%%%%%%% COMMANDES %%%%%%%%%%%%%%%%%%%%%%%%%%%

\newcommand{\itemb}{\item[$\bullet$]}

%% Les ensembles de modelisation

\newcommand{\Fd}{\mathbb{F}}
\newcommand{\Znat}{\mathbb{Z}}
\newcommand{\Nnat}{\mathbb{N}}

%%%%%%%%%%%%%%% ENVIRONEMENTS %%%%%%%%%%%%%%%%%%%%%%%


\theoremstyle{plain}
\newtheorem{lemme}{Lemme}
\newtheorem{algorithme}{ Algorithme }
\newtheorem{corolaire}{Corollaire}
\newtheorem{propriete}{Propri\'et\'e}
\newtheorem{proposition}{Proposition}
\newtheorem{theoreme} {Th\'eor\`eme}
\newtheorem{definition}{D\'efinition}

\theoremstyle{definition}
\newtheorem{exercice}{Exercice}
%\theoremstyle{remark}
\newtheorem{remarque}{Remarque}
\newtheorem{exemple}{Exemple}


%%%%%%%%%%%%%% DOCUMENTS %%%%%%%%%%%%%%%%%%%%%%%%%%%%


\begin{document}

\hbox{\raisebox{0.4em}{\vrule depth 0pt height 0.5pt width \textwidth}}
\noindent \textbf{Universit� de Perpignan} \hfill  L3 Informatique \\
 \hfill C. Legrand \hfill  Ann�e \the\year \\
\smallskip
\begin{center}
{
  {\scshape \large TD2 S�curit� \\ Programmes malveillants}  \\
 %G�n�ralit�
}



\end{center}
\hbox{\raisebox{0.4em}{\vrule depth 0pt height 0.9pt width \textwidth}}
%\hline

\bigskip

\bigskip


\begin{exercice}
\begin{enumerate}
\item Dans quelle mesure les vers sont-ils plus dangereux que les virus?
\item Certains vers qui se propagent sur Internet ne provoquent aucun dommage sur les machines atteintes. Pourquoi sont-ils cependant nuisibles?
\item Rappeler deux techniques qui permettent � des virus de ne pas �tre d�tect�s par des logiciels antivirus.
\end{enumerate}
\end{exercice}



\begin{exercice}
On consid�re le programme (en langage C) suivant
\begin{verbatim}
//... les includes
int main(int argc, char **argv)
{
  if(argc >= 1)
    {	
       system("cp argv[0] argv[1]");
    }
}
\end{verbatim}
est ce un virus? 
\end{exercice}


\begin{exercice}
On consid�re dans cet exercice une variante du ver W32/Beagle. Ce ver se pr�sente sous forme d'un courrier �lectronique poss�dant un fichier joint, qui est � la fois compress� et chiffr�. Le mot de passe pour d�chiffrer le fichier est contenu dans le corps du message. Si la victime ex�cute le fichier obtenu apr�s d�compression avec le mot de passe fourni (qui est un fichier avec une extension \texttt{.exe}), alors le ver se propage en choisissant la prochaine victime dans le carnet d'adresses de la victime courante. 

Pourquoi le fichier compress� est-il chiffr� puisque le mot de passe est fourni dans le message?

\end{exercice}


%\begin{exercice}
%\'Ecrire un virus par �crasement de code (en Langage C, ou en un langage de script). Imaginer une/des strat�gie pour que la taille du fichier infect� ne change pas.
%\end{exercice}


\begin{exercice}
Voici une descrition d'un logiciel malveillant \texttt{Kork} 
\begin{verbatim}
Kork
ALIAS:	Linux/Kork, Unix/Kork, 

Kork  uses the known vulnerability in lpd service to propagate from a vulnerable 
Linux system to another. This service is part of the default installation of Red 
Hat Linux 7.0.

VARIANT:	Kork.A
If Kork.A finds a vulnerable host, it first creates two users, "kork" and "kork2", 
to the system without a password. "kork2" user has a root priviledge. Kork also 
adds an open shell to port 666.

Next it attempts to download a trojanized login and the main part of the trojan 
from a web site. Since April 26th, 2001, neither of these files are available, so 
it cannot replicate any further. However, already infected machines are able to 
compromise other vulnerable machines by adding an open shell and users to the 
system.

If the download is completed, Kork install an infected "/bin/login" and "/bin/ps" 
to the system. It attemps to send sensitive system data propably to the author of 
Kork. Original "/bin/ps" is copied to "/usr/bin/.ps" and the original "/bin/login" 
is copied to "/bin/.login". Finally Kork starts to scan random Class-B subnets for 
vulnerable hosts. 
\end{verbatim}

\begin{enumerate}
\item Dire de quel type de programme malveillant est Kork (bombe logique, cheval de Troie, etc.) ?
\item Quels sont les cibles potentielles du virus? Evaluer bri�vement son impact.
\end{enumerate}
\end{exercice}

\begin{exercice}
On consid�re un rapport concernant un ex�cutable malveillant nomm� \texttt{Staog}.
\begin{verbatim}
NAME: Staog
SIZE: 4744

Staog spreads only under Linux operating system, infecting
Elf-style executables.Its first appearance was  in the fall of 1996, 

Staog is written in assembler. It attempts to stay resident and
infect binaries as they are executed by any user. Stoag tries to
subvert root access via three known vulnerabilities (mount buffer
overflow, tip buffer overflow and one suidperl bug).

Staog contains several text strings, including:

Staog by Quantum / VLAD
/dev/kmemx/etc/mtab~
/sbin/mount
/tmp/t.dip
/bin/sh
/sbin/dip /tmp/t.dip
chatkey
/tmp/hs
#!/bin/sh\nchmod 666 /dev/kmem\n/tmp/hs
#!/usr/bin/suidperl -U\n$ENV{PATH}=\"/bin:/usr/bin\";
\n$>=0;$<=0;\nexec(\"chmod 666 /dev/kmem\");\n


Staog can be detected by searching all binaries for the following hex
search string:

215B31C966B9FF0131C0884309884314B00FCD80

Staog is not known to be in the wild at the time of this writing
(February 1997).
\end{verbatim}

\begin{enumerate}
\item Dire de quel type de programme malveillant est Staog (bombe logique, cheval de troie, etc.) ?
\item Quels sont ses cibles potentielles et son mode d'action, et �valuer bri�vement son impact.
\item Quelle est sa stat�gie de furtivit� et sa faiblesse ?
\end{enumerate}

\end{exercice}

\end{document}
